\chapter{提案内容 or 提案手法}
  本章では、本研究の提案内容について述べる． 
  \section{全体概要}
  本研究では，3 つの研究課題を設定した．1 つ目にYouTube ライブ配信コメントの特徴とXのポストとの比較．2 つ目が YouTube に投稿されたコメントを長期的に収集し分析を可能とする基盤の提案．3 つ目に，地震発生からコメントのピーク検出，感情分析モデルを適用し比較検討など，収集したコメントに対して分析を行い，特徴を明らかにする．
  
  研究課題の1つ目でYouTubeライブ配信コメントの仕様やユーザ属性の違いなどを理解した上で，研究課題2である、ライブ配信に投稿されたコメントを長期的に収集し分析を可能とする基盤の提案を行う．最終的に，研究課題3において収集されたコメントを対象に、投稿されたコメント数の時系列変化やピーク検出、感情分析モデルの適用を行い、SNSから災害時のYouTubeライブ配信コメントの特徴を明らかにする．
  
  \subsection{XとYouTubeライブ配信コメントの仕様についての比較（理由）}
    % \begin{itemize}
    %   \item XとYouTubeライブ配信コメントの仕様についての比較（理由）
    %   \item 収集量の仕様比較
    %   \item これらにより、Xよりも大量のデータを収集できる余地があると示す
    % \end{itemize}
  \subsection{YouTubeライブ配信コメントの収集について}
  
  \subsection{収集した情報を元に分析する（複数、実際の記述は順不同で説明に応じて前後）}
    \begin{itemize}
      \item Xでも行われている分析をある程度やりたい（細かい方法ではなくこれらからわかることや応用先を書く）
      
      1. 地震からコメントのピークや立ち上がり時間の分析
      
      2. コメントのクラスタリング
      
      3. 感情分析
      
      4. コメントユーザの分類
      
      5. LLMを利用した分析
      
      6. など...

    \end{itemize}

\section{XとYouTubeライブ配信コメントの仕様についての比較}
\begin{itemize}
    \item 比較する仕様や内容の枠を設定
    \item など…
\end{itemize}

\section{YouTubeライブ配信コメントの収集について}
\begin{itemize}
    \item システム概要
    \item 収集
    \item 蓄積
    \item 分析部分
    \item 可視化
\end{itemize}

\section{収集した情報の分析（細かい方法などを記述）}
      1. 地震からコメントのピークや立ち上がり時間の分析
      
      2. コメントのクラスタリング
      
      3. 感情分析
      
      4. コメントユーザの分類
      
      5. LLMを利用した分析
      
      6. など...